%2025 - 9b

{\bf Chapitre 6 : Métrique et géodésiques (suite)}
"Compatibilité" entre métrique et dérivée covariante

$(M, \mc{O}, \mc{A}, g, ^{LC}\nabla)$
\[
\tag{1}(\nabla_{\frac{\partial}{\partial x^a}}g)_{bc} =
\frac{\partial}{\partial x^a} g_{bc}-\Gamma^i_{ca}\,g_{ib}-\Gamma^j_{ba}\,g_{cj}
\]
\[
\tag{2}(\nabla_{\frac{\partial}{\partial x^b}}g)_{ca} =
\frac{\partial}{\partial x^b} g_{ca}-\Gamma^i_{ab}\,g_{ic}-\Gamma^j_{cb}\,g_{aj}
\]
\[
\tag{3}(\nabla_{\frac{\partial}{\partial x^c}}g)_{ab} =
\frac{\partial}{\partial x^c} g_{ab}-\Gamma^i_{bc}\,g_{ia}-\Gamma^j_{ac}\,g_{bj}
\]

(1)-(2)-(3)


\[
0=\frac{\partial}{\partial x^a} g_{bc}-\frac{\partial}{\partial x^b} g_{ca}
-\frac{\partial}{\partial x^c} g_{ab}-2\Gamma^i_{cb}g_{ia}
\]

\[
g_{ia}\Gamma^i_{cb}=\frac{1}{2}\left(\frac{\partial}{\partial x^b} g_{ca} +
\frac{\partial}{\partial x^c} g_{ba}-\frac{\partial}{\partial x^a} g_{bc}\right)
\]

\[
g^{ka}_{jk}
\]

Connexion de Levi-Civita

dérivé covariante, "commute" avec les "montées/descentes" d'indice

\[
\nabla_{\frac{\partial}{\partial x^i}}(Bra(X)) =
Bra(\nabla_{\frac{\partial}{\partial x^i}}X)
\]

\[
Bra(\nabla_{\frac{\partial}{\partial x^i}}X) =
g_{kj}\nabla_{\frac{\partial}{\partial x^i}}^k
\]

\[
X = X^i\frac{\partial}{\partial x^i} \qquad Bra(X) = X_jdx^j \quad
o\grave u  \quad X_j = g_{ij}X^i
\]
\[
\nabla_Y(Bra(X))=Bra(\nabla_YX) \qquad \nabla_{\frac{\partial}{\partial x^i}}X
\]

Tenseur de courbure de Ricci
\[
(\nabla_Xg)_j=0
\]


Courbure scalaire de Ricci

\[
\begin{array}{ c c l }
 \mc{R}^i_{jkl} & \to & \mc{R}^i_{jil} \equiv R_{jl} \\
  Riemann &  & Ricci \\
    & \xrightarrow[\text{g}]{} & g^{jk}R_{jl} = R^k_l \to R^k_k=R \\
\end{array}
\ \ 
\]
\[
g,\ (1,3) \quad \to \quad ^\mc{U}\Gamma \quad symboles\ de\ Christoffel
\]
\[
\mc{R}^i_{jkl}=\partial\Gamma-\partial\Gamma+\Gamma\Gamma-\Gamma\Gamma
\]
\[
\frac{\partial g}{\partial x^i} \qquad (\nabla _X g)_j=0
\]
\[
\mc{R}^i_{jkl} \to \frac{\partial ^2g}{\partial x^i\partial x^j}
\]

Tenseur d'Einstein (et équations d'Einstein)
\[
\boxed{ G_{\mu\nu} = R_{\mu\nu}-\frac{1}{2}Rg_{\mu\nu} }
\qquad G_{\mu\nu} = \kappa T_{\mu\nu}
\]
$T_{\mu\nu}$ est le tenseur énergie impulsion
\[
S=\int \underbrace{\mc{R}\sqrt{-g}}_{\displaystyle \mc{L}_{espace-temps}}\ d
 + \mc{L}_{mati\grave{e}re}
\]
$\mc{L}_{mati\grave{e}re}$ correspondant à $T_{\mu\nu}$


%Remarque :
%gravitation newtonienne et espace-temps courbe
\section{Gravitation newtonienne}
\[
m_{inerte}\frac{d\vec{v}}{dt}=\vec{F}_{Grav}=m_{Grav}\vec{g}
\]
\[
\ddot{x}^i_\gamma-\frac{F}{m}^i(\gamma)=0
\]
\[
\ddot{x}^i_\gamma-\frac{F}{m}^i(\gamma)
\times\underset{\substack{\downarrow \\ \frac{dt}{dt}}}{1}
\times\underset{\substack{\downarrow \\ \frac{dt}{dt}}}{1}=0
\]

\[
\ddot{x}^i_\gamma-\frac{F}{m}^i\dot{x}^0\dot{x}^0=0
\quad (x^0=t,\ x^1,\ x^2,\ x^3)
\]
\[
\Gamma^i_{jk} = 0 \quad sauf \quad \Gamma^i_{00}=-\frac{F}{m}^i
\]
\[
\frac{\vec{F}}{m}=\vec{g}=-\vec{grad}\phi
\qquad \Delta\phi=4\pi G\rho
\]
\[
R_{00}=4\pi G T_{\mu\nu}
\]
\[
g^{\mu\nu}G_{\mu\nu}=R-\frac{1}{2}Rg^{\mu\nu}g_{\mu\nu}
 = 8\pi Gg^{\mu\nu}T_{\mu\nu}
\]
\[
R_{\mu\nu}=8\pi G(T_{\mu\nu}-\frac{1}{2}Tg_{\mu\nu})
\]



